\chapter*{Course Guide}
\addcontentsline{toc}{chapter}{Course Guide}
\label{front:syllabus}

\section*{Aim of the Course}
\addcontentsline{toc}{section}{Aim of the Course}

This twelve-week course gives a route through the book for one two-hour lecture
each week.  The goal is not to survey elliptic curves broadly, but to understand
one central identification in detail:
\[
  \C/\Lambda \;\cong\; \Pic^0(E) \;\cong\; E(\C).
\]
Here the isomorphisms are isomorphisms of complex Lie groups.  The course ends
by proving that the chord--tangent law on a smooth plane cubic is addition on
the torus transported by
\[
  z\longmapsto[\wp(z):\wp'(z):1].
\]

\section*{Prerequisites}
\addcontentsline{toc}{section}{Prerequisites}

The course assumes a first course in complex analysis.  The following table
separates what is used immediately from what the text develops during the
course.

\begin{center}
\small
\renewcommand{\arraystretch}{1.18}
\begin{tabular}{@{}L{0.24\linewidth}L{0.64\linewidth}@{}}
\textbf{Assumed} &
Holomorphic functions, contour integration, residues, isolated singularities,
and basic topology of subsets of \(\C\).\\
\textbf{Reviewed as needed} &
Differential forms, Stokes' theorem, meromorphic functions on compact surfaces,
and zero-pole counting.\\
\textbf{Developed carefully} &
Compact Riemann surfaces, divisors, Picard groups, Jacobians,
genus-one Riemann--Roch, and the Weierstrass embedding.
\end{tabular}
\end{center}

\section*{Structure of the Course}
\addcontentsline{toc}{section}{Structure of the Course}

\begin{center}
\small
\renewcommand{\arraystretch}{1.22}
\begin{tabular}{@{}L{0.11\linewidth}L{0.23\linewidth}L{0.27\linewidth}L{0.27\linewidth}@{}}
\textbf{Weeks} & \textbf{Theme} & \textbf{Main objects} & \textbf{Question}\\
\hline
1--3 &
Analytic foundations &
Forms, residues, compact surfaces, lattices &
How does complex analysis produce a compact genus-one group?\\
4--6 &
Elliptic functions &
Periodic meromorphic functions, \(\wp\), principal divisors &
How do functions on \(\C/\Lambda\) remember the group law?\\
7--9 &
Divisors and Jacobians &
\(L(D)\), \(\Pic^0(E)\), Abel--Jacobi, \([\wp:\wp':1]\) &
Why does the analytic torus become a smooth plane cubic?\\
10--12 &
The cubic group law &
Lines, tangents, Abel's theorem, \(j\)-invariant &
Why is chord--tangent addition the same operation as torus addition?
\end{tabular}
\end{center}

\section*{Weekly Lecture Plan}
\addcontentsline{toc}{section}{Weekly Lecture Plan}

\begin{center}
\footnotesize
\renewcommand{\arraystretch}{1.16}
\begin{tabular}{@{}L{0.06\linewidth}L{0.31\linewidth}L{0.15\linewidth}L{0.36\linewidth}@{}}
\textbf{Week} & \textbf{Lecture focus} & \textbf{Reading} & \textbf{Core result}\\
\hline
1 &
Forms, residues, and global obstructions &
Ch. 1; App. A &
Residues turn local singular data into a global constraint.\\
2 &
Compact Riemann surfaces and genus &
Ch. 2 &
Compactness changes complex analysis from local function theory into global
geometry.\\
3 &
Lattices and complex tori &
Ch. 3 &
\(\C/\Lambda\) is a compact genus-one Riemann surface and a complex Lie group.\\
4 &
Elliptic functions &
Ch. 4 &
Zeros, poles, and residues of elliptic functions balance on a fundamental
parallelogram.\\
5 &
The Weierstrass \(\wp\)-function &
Ch. 5 &
\(\wp\) and \(\wp'\) satisfy
\(\wp'^2=4\wp^3-g_2\wp-g_3\).\\
6 &
Divisors and principal divisors &
Ch. 6 &
Zeros and poles become degree-zero divisor classes modulo principal divisors.\\
7 &
Riemann--Roch in genus one &
Ch. 7 &
\(\ell(nO)=n\) for \(n\geq 1\), giving the functions needed for the cubic map.\\
8 &
Abel's theorem and the Jacobian &
Ch. 8 &
The map \(P\mapsto[P]-[O]\) identifies \(E\) with \(\Pic^0(E)\) as complex Lie
groups.\\
9 &
The Weierstrass embedding &
Ch. 9 &
\(1,\wp,\wp'\) embed the torus into \(\CP^2\) as a smooth cubic.\\
10 &
The plane cubic group law &
Ch. 10 &
Three collinear points on the cubic have sum \(O\).\\
11 &
Transport of addition &
Ch. 11; Ch. 14 &
The chord--tangent operation is torus addition transported by
\([\wp:\wp':1]\).\\
12 &
Synthesis and further directions &
Chs. 12--14 &
Cohomology, moduli, the \(j\)-invariant, and finite-field questions extend the
same genus-one structure.
\end{tabular}
\end{center}

\section*{Exercises and Written Work}
\addcontentsline{toc}{section}{Exercises and Written Work}

Exercises are meant to force translation between the analytic, divisor-theoretic,
and geometric languages.  They should be attempted before consulting the
solutions appendix.

\begin{center}
\small
\renewcommand{\arraystretch}{1.2}
\begin{tabular}{@{}L{0.22\linewidth}L{0.48\linewidth}L{0.18\linewidth}@{}}
\textbf{Work} & \textbf{Main task} & \textbf{Due}\\
\hline
Problem Set 1 &
Residues, compactness, genus, and quotient tori &
After Week 3\\
Problem Set 2 &
Elliptic functions, \(\wp\), and principal divisors &
After Week 6\\
Problem Set 3 &
Riemann--Roch, Abel--Jacobi, and the Weierstrass embedding &
After Week 9\\
Problem Set 4 &
Line divisors, tangents, and transported addition &
After Week 11\\
Final synthesis &
A concise proof narrative of
\(\C/\Lambda\cong\Pic^0(E)\cong E(\C)\) as complex Lie groups &
Week 12
\end{tabular}
\end{center}

\begin{keyidea}{Final target}
By the end of the course, the student should be able to read the equation of
the Weierstrass cubic, the divisor class \([P]-[O]\), and the point
\(z\in\C/\Lambda\) as three notations for the same genus-one group element.
\end{keyidea}
